\addchap{Thesen}
\begin{enumerate}
  \item Ungelabelte Bilder können Deep Metric Learning auf eingefrorenen DINOv2-Merkmalen verbessern, wenn wenige Labels vorliegen, zusätzliche echte Labels die supervised Referenz noch deutlich verbessern und die ungelabelten Bilder aus den gelabelten Klassen stammen. Bei den Budgets der Hyperparametersuche gewinnt Label Propagation auf Cars196 und In-Shop in jedem gepaarten Vergleich, um 2,2 bzw. 3,2 Prozentpunkte. Mit der übertragenen Konfiguration erreicht sie auf Cars196 mit vier gelabelten Trainingsbildern je Klasse die Referenz mit sechzehn.
  \item Der Nutzen ungelabelter Bilder sinkt mit dem Label-Budget. Bei 32 gelabelten Bildern je Klasse (Cars196) liegt er im Mittel unter einem Prozentpunkt, bei 60\,\% gelabeltem Pool (In-Shop) ist er verschwunden.
  \item Sind die vortrainierten Merkmale bereits nahe am Erreichbaren, wie auf \ac{CUB}, bleibt selbst ein supervised trainierter Head mit wenigen Labels als Einzelmodell unter den unveränderten Merkmalen. Ungelabelte Bilder bringen dort gegenüber der supervised Referenz nur kleine Gewinne, vor allem im Ensemble.
  \item Die für die semi-supervised Verfahren anders gewählten Hyperparameter erklären die Gewinne nicht. Dass die ungelabelten Bilder auch mit den Hyperparametern der supervised Referenz helfen, gilt aber vor allem für Label Propagation. Je nach Verfahren und Verlustfunktion hebt schon der semi-supervised Term auf den gelabelten Bildern allein die Leistung, bei der Laplace-Regularisierung auf Cars196 mit allen fünf Verlustfunktionen und bei Hoffer auf In-Shop mit vier.
  \item Label Propagation und Laplace-Regularisierung wirken über die aus den Daten berechnete Nachbarschaft, denn ein zufällig umverdrahteter Graph schadet stark.
  \item Für Pseudo-Labels im Metric Learning ist die Fehlerrate allein kein ausreichendes Maß. Wie stark falsche Labels schaden, hängt auch von der Verlustfunktion und von den betroffenen Bildern ab. Mit einem Seed liegt Label Propagation auf Cars196 mit drei der fünf und auf In-Shop mit allen fünf Verlustfunktionen über Läufen mit zufälligen Fehlern gleicher mittlerer Häufigkeit.
  \item Ob ein Verfahren als besser gilt, hängt davon ab, ob ein Einzelmodell oder ein Ensemble der Fold-Modelle bewertet wird. Beide Sichtweisen müssen deshalb getrennt berichtet werden.
  \item Stammen die ungelabelten Bilder aus Klassen ohne Labels, schadet Label Propagation in der hier verwendeten Form, die stets einen Teil der Bilder übernimmt. Die Laplace-Regularisierung hilft bei wenigen gelabelten Klassen im Ensemble weiter, als Einzelmodell aber nur mit Circle.
\end{enumerate}
